\chapter{Blood Groups}
\index{Blood Groups}

\section*{Definition and Overview}
Blood groups (blood types) are classifications of blood based on the presence or absence of inherited markers called antigens on the surface of red blood cells. The most clinically important systems are the ABO system, which divides blood into groups A, B, AB, and O, and the Rhesus (Rh) system, which marks blood as Rh positive or Rh negative. Blood group is inherited and normally has no effect on health, but it is critical for blood transfusion and for pregnancy. If a person receives blood of an incompatible ABO group, naturally occurring antibodies can attack the transfused cells and cause a serious, potentially life-threatening haemolytic transfusion reaction. Rh status matters when an Rh negative mother carries an Rh positive baby, because her immune system may produce antibodies that cross the placenta and harm the baby's red cells.

\section*{Epidemiology}
Blood group distribution varies between populations. Worldwide, group O is the most common (roughly 4 in 10 people), followed by groups A, B, and AB, which is the rarest. About 8 in 10 people are Rh positive. In many countries, O positive is the most common group and AB negative the rarest. Because a person can only safely receive certain groups, blood services must maintain balanced stocks of all groups, and the distribution of groups within a community influences donor recruitment, transfusion planning, and the availability of compatible blood for people with rare types.

\section*{Aetiology and Pathophysiology}
Blood group is determined by genes inherited from both parents. The ABO group reflects which sugars (antigens) are present on the red cell surface: group A cells carry A antigen, group B cells carry B antigen, group AB cells carry both, and group O cells carry neither. A person naturally produces antibodies against the A or B antigens they do not possess, so group A blood makes anti-B, group B makes anti-A, and group O makes both. These naturally occurring antibodies appear in early life and are the reason a mismatched ABO transfusion causes an immediate immune attack on the donor cells.

The Rh system operates differently. Most people carry the Rh D antigen and are Rh positive; those who lack it are Rh negative. Rh negative people do not naturally produce anti-D antibodies, but they can develop them after exposure to Rh positive blood through a transfusion or during pregnancy with an Rh positive baby. If an Rh negative woman develops anti-D antibodies and later carries another Rh positive baby, the antibodies can cross the placenta and destroy the fetal red cells, causing haemolytic disease of the newborn.

\begin{itemize}
    \item \textbf{ABO antigens:} A, B, AB, or O, determined by inherited genes
    \item \textbf{Rh D antigen:} present (Rh positive) or absent (Rh negative)
    \item \textbf{Naturally occurring antibodies:} anti-A and anti-B, present from early life
    \item \textbf{Immune antibodies:} anti-D develops only after exposure to Rh positive blood
    \item \textbf{Rare other systems:} Kell, Duffy, Kidd, and MNS matter mainly for repeated transfusions and for some cases of haemolytic disease
\end{itemize}

\section*{Clinical Features}
A person's blood group causes no symptoms, and most people discover theirs only through testing, blood donation, or pregnancy screening. Problems arise only from blood group incompatibility:

\begin{itemize}
    \item \textbf{Transfusion reaction:} fever, chills, pain in the back or flank, dark urine, breathlessness, and a fall in blood pressure during or shortly after an incompatible transfusion
    \item \textbf{Haemolytic disease of the newborn:} jaundice, pallor, and anaemia in a baby whose mother's antibodies destroy its red cells
    \item \textbf{Rh sensitisation:} no symptoms at the time, but it places future pregnancies at risk
    \item \textbf{Difficulty finding compatible blood:} people with rare groups may experience delays when transfusion is needed
\end{itemize}

\section*{Differential Diagnosis}
Reactions and findings caused by blood group incompatibility must be distinguished from other conditions:

\begin{itemize}
    \item Other transfusion reactions: allergic reactions, febrile non-haemolytic reactions, transfusion-related acute lung injury, or transfusion-transmitted infection
    \item Newborn jaundice from other causes, such as prematurity, sepsis, ABO incompatibility, or enzyme deficiency
    \item Anaemia in pregnancy from iron deficiency or other causes rather than Rh disease
    \item Other causes of acute shock and kidney injury occurring during a transfusion
\end{itemize}

\section*{When to Seek Medical Care}
Blood group testing itself requires no medical care, but any reaction during a transfusion must be reported to the attending staff immediately, since rapid recognition saves lives. Pregnant women who are Rh negative should ensure they receive routine antenatal blood group and antibody testing, and anti-D prophylaxis as recommended. Seek medical advice promptly if a baby becomes deeply jaundiced in the first days of life, appears pale, or is unusually sleepy and feeds poorly.

\textbf{Contact your local emergency services for:}
\begin{itemize}
    \item Fever, chills, or difficulty breathing during a blood transfusion
    \item Pain in the back, chest, or abdomen during a transfusion
    \item Signs of a severe allergic reaction, such as swelling of the face or throat, or widespread hives
    \item Collapse, loss of consciousness, or a rapid drop in blood pressure
\end{itemize}

\section*{Diagnosis and Investigations}
\begin{itemize}
    \item \textbf{Blood group typing:} determines ABO group and Rh type from a blood sample
    \item \textbf{Antibody screen:} checks for unexpected red cell antibodies, including anti-D
    \item \textbf{Crossmatch:} before transfusion, donor and recipient blood are tested together for compatibility
    \item \textbf{Antenatal testing:} maternal blood group and antibody screening in early pregnancy, repeated in Rh negative women
    \item \textbf{Baby's blood test:} cord blood or heel-prick samples can confirm a newborn's group and check for anaemia and bilirubin levels
\end{itemize}

\section*{Management}
\begin{itemize}
    \item \textbf{Compatible transfusion:} patients always receive blood matched for ABO and Rh, with a crossmatch performed before transfusion
    \item \textbf{Emergency O negative blood:} used when there is no time to determine a patient's group
    \item \textbf{Anti-D prophylaxis:} Rh negative women receive anti-D immunoglobulin after events that could cause sensitisation, and sometimes in late pregnancy, to protect future pregnancies
    \item \textbf{Monitoring during transfusion:} vital signs are checked and the patient observed during and after the transfusion
    \item \textbf{Treatment of reactions:} stopping the transfusion immediately and providing urgent supportive care, including fluids and treatment for shock
    \item \textbf{Treatment of affected newborns:} phototherapy for jaundice, or exchange transfusion for severe haemolytic disease
    \item \textbf{Registration of rare donors:} people with rare blood groups can be registered with blood services so compatible blood can be located if needed
\end{itemize}

\section*{Complications}
\begin{itemize}
    \item Acute haemolytic transfusion reaction from an ABO mismatch, which can cause kidney failure, shock, and disseminated intravascular coagulation
    \item Delayed haemolytic reactions occurring days after a transfusion
    \item Rh disease in pregnancy, which can cause severe fetal anaemia, heart failure, and hydrops fetalis
    \item Kernicterus (brain damage from severe newborn jaundice) if haemolytic disease is untreated
    \item Sensitisation, which makes future transfusions and pregnancies more difficult to manage safely
\end{itemize}

\section*{Prognosis and Outcomes}
With modern blood typing, crossmatching, and anti-D prophylaxis, complications of blood group incompatibility are rare, and transfusions are very safe when the blood group is correctly matched. Rh disease of the newborn has become uncommon since anti-D prophylaxis became routine, and affected babies usually recover fully with treatment. People with rare blood groups may require more planning, but registered donor networks and national rare-donor programmes help meet their needs.

\section*{Prevention and Self-Care}
\begin{itemize}
    \item Know your blood group and, if relevant, inform healthcare providers before transfusions or surgery
    \item Wear medical-alert identification if you have a rare blood group or unusual antibodies
    \item Attend routine antenatal blood tests if pregnant, so that Rh problems can be prevented
    \item Receive anti-D prophylaxis on schedule if you are Rh negative and pregnant or after any sensitising event
    \item Donate blood regularly if you can -- every group is needed, and rare groups are especially valuable
    \item Seek early medical attention if a baby develops jaundice in the first days of life
\end{itemize}

\section*{Sources and Further Reading}
The following sources provide reliable, up-to-date information on blood groups, transfusion safety, and pregnancy screening:

\begin{itemize}
    \item NHS. \textit{NHS guidance on blood groups and transfusion.} https://www.nhs.uk/conditions/
    \item Mayo Clinic. \textit{Information on blood types and blood donation.} https://www.mayoclinic.org/
    \item World Health Organization. \textit{Blood transfusion safety resources.} https://www.who.int/
    \item MedlinePlus. \textit{Blood typing and crossmatch patient information.} https://medlineplus.gov/
    \item MSD Manual. \textit{Overview of blood groups and transfusion reactions.} https://www.msdmanuals.com/
    \item Australian Red Cross Lifeblood. \textit{Information about blood types and donation.} https://www.lifeblood.com.au/
\end{itemize}
